\chapter{Frequently Asked Questions and Pitfalls}
\label{app:faq}

This appendix summarizes problems and questions frequently encountered by readers while studying Volume 1. Each item deals with specific situations that may occur during actual implementation.

\section{environmental related}

\subsection{Q: PyTorch After installation CUDAis not recognized}

\textbf{A}: Check two things. First, make sure you have the NVIDIA drivers installed (run\code{nvidia-smi}). Second, whether the PyTorch version and CUDA version are compatible. For example, if you install PyTorch for CUDA 12.1 on a CUDA 11.8 system, it will not be recognized.

\begin{lstlisting}[style=bashstyle]
# CUDA Check version
nvidia-smi | grep "CUDA Version"

# PyTorch Reinstall (CUDA 12.For 1)
pip uninstall torch
pip install torch --index-url https://download.pytorch.org/whl/cu121
\end{lstlisting}

\subsection{Q: Learning is too slow}

\textbf{A}: Check the following:
\begin{enumerate}
\item GPU is being used (set\code{device="cuda"}).
\item Is the batch size too small (low GPU utilization)?
\item Is data loading the bottleneck (increases\code{num\_workers}).
\item Are you using mixed precision (BF16)?
\end{enumerate}

\section{Tokenizer related}

\subsection{Q: BPEThe tokenizer learned with ``hi''It tokenizes strangely.}

\textbf{A}: Because Korean was not sufficiently included in the corpus. BPE is frequency-based, so rare characters are broken down into bytes. way out:
\begin{itemize}
\item Increase the Korean corpus ratio.
\item Increase vocabulary size (about 20,000 Korean syllables + subwords).
\item Use a professional tokenizer like SentencePiece.
\end{itemize}

\subsection{Q: encode/decode Round trip fails}

\textbf{A}: Check space handling. The `` '' character in BPE may not be converted to a space when decoding. In our implementation, the\code{\_postprocess}method handles this:

\begin{lstlisting}[language=Python]
def _postprocess(self, text):
    text = text.replace("▁", " ")
    return text.strip()
\end{lstlisting}

\subsection{Q: special token IDWhat happens if changes}

\textbf{A}: Fatal. The model was trained assuming\code{<pad>=0, <bos>=1, <eos>=2, <unk>=3}. If the special token ID changes while retraining the tokenizer, the model weights are no longer valid. Always fix special token IDs.

\section{model related}

\subsection{Q: Attention output NaNThis comes out}

\textbf{A}: Three causes are common:
\begin{enumerate}
\item\textbf{Softmax saturation}: Logit is too large, so softmax is saturated at 0/1. Check scaling ($\sqrt{d\_k}$).
\item\textbf{Learning rate too large}: Gradient explosion. Try reducing the learning rate to 1/10.
\item\textbf{FP16 overflow}: The activation value exceeds the FP16 range ($\pm 65504$). Switch to BF16.
\end{enumerate}

\begin{lstlisting}[language=Python]
# NaN check
if torch.isnan(loss):
    print("NaN detected! Check learning rate and gradients.")
    print(f"Max grad: {max(p.grad.abs().max() for p in model.parameters() if p.grad is not None)}")
\end{lstlisting}

\subsection{Q: The model repeats the same word}

\textbf{A}: Check your sampling strategy. Greedy is prone to falling into repetitive loops. Switch to Top-P (P=0.9) or set Temperature to 0.7$\sim$0.9. Adding a repetition penalty is also effective.

\begin{lstlisting}[language=Python]
# Repetition penalty (simplified)
def apply_repetition_penalty(logits, generated_ids, penalty=1.2):
    for token_id in set(generated_ids.tolist()):
        logits[token_id] /= penalty
    return logits
\end{lstlisting}

\subsection{Q: An error occurs if the context length is exceeded.}

\textbf{A}: If input exceeding\code{context\_length}of the model is entered, an error occurs at the embedding stage.  Notice the cut to the sliding window in the\code{generate}function:

\begin{lstlisting}[language=Python]
if input_ids.shape[1] > context_length:
    x = input_ids[:, -context_length:]  # Keep only the most recent
\end{lstlisting}

\section{learning related}

\subsection{Q: Losses do not decrease}

\textbf{A}: Check the following:
\begin{enumerate}
\item The learning rate is either too small (1e-6 or less) or too large (1e-2 or more).
\item Too little data (at least a few thousand tokens required).
\item Bad optimizer (check AdamW's beta, eps).
\item Gradient does not flow (check requires\_grad).
\end{enumerate}

\begin{lstlisting}[language=Python]
# Check gradient flow
for name, p in model.named_parameters():
    if p.requires_grad and p.grad is None:
        print(f"No gradient for {name}")
    elif p.grad is not None and p.grad.abs().max() < 1e-10:
        print(f"Vanishing gradient for {name}")
\end{lstlisting}

\subsection{Q: Loss occurs in the early stages of learning.}

\textbf{A}: Possibility of insufficient warmup or incorrect initialization. Try this:
\begin{itemize}
\item Increase Warmup Steps (100$\sim$1000).
\item Check for gradient clipping (\code{grad\_clip=1.0}).
\item Reduce initialization standard deviation (0.02$\to$0.01).
\end{itemize}

\subsection{Q: Loading a checkpoint gives different results}

\textbf{A}: Two causes:
\begin{enumerate}
\item\textbf{Seed not set}: Check whether\code{set\_seed(42)}has been called.
\item\textbf{dropout}: Check whether the model has been switched to\code{eval()}mode.
\item\textbf{CuDNN indeterminism}:\code{torch.backends.cudnn.deterministic = True}.
\end{enumerate}

\section{Distributed Learning (Preview of Volume 2)}

\subsection{Q: DDPIf you learn with , the speed is rather slow.}

\textbf{A}: Possible communication bottleneck. Check the following:
\begin{itemize}
\item The batch size is too small, so the computation/communication ratio is low.
\item Use slower communications (PCIe, etc.) rather than NVLink.
\item Gradient synchronization occurs at every step (reduced to gradient accumulation).
\end{itemize}

\subsection{Q: FSDPat OOMThis is me}

\textbf{A}: Not enough activation memory. Turn on checkpointing or reduce batch size. CPU offload is also considered.

\section{Debugging Tips}

\subsection{Verify step by step}

Complex bugs are broken down into small units and verified. Testing only the tokenizer, testing only the attention, etc.

\begin{lstlisting}[language=Python]
# Tokenizer standalone test
tok = BPETokenizer()
tok.train(["hello world"], vocab_size=50)
print(tok.encode("hello"))
print(tok.decode(tok.encode("hello")))

# Attention stand-alone test
attn = MultiHeadAttention(config)
x = torch.randn(1, 4, 32)
out = attn(x)
print(out.shape, out.dtype, torch.isnan(out).any())
\end{lstlisting}

\subsection{shapeAlways print}

80% of bugs are shape mismatches. The habit of printing shapes at every step.

\begin{lstlisting}[language=Python]
print(f"input shape: {x.shape}")
print(f"after embedding: {emb.shape}")
print(f"after attention: {attn_out.shape}")
\end{lstlisting}

\subsection{Start with small input}

Start with batch 1, sequence 4, check operation, then increase size. Debugging on large inputs can be confusing because there is too much output.

\section{Performance Optimization Tips}

\subsection{Use mixed precision}

\begin{lstlisting}[language=Python]
from torch.cuda.amp import autocast, GradScaler

scaler = GradScaler()
with autocast(dtype=torch.bfloat16):
    logits = model(inputs)
    loss = criterion(logits, targets)
scaler.scale(loss).backward()
scaler.step(optimizer)
scaler.update()
\end{lstlisting}

Unlike FP16, BF16 does not require a scaler. Recommended for A100 and above.

\subsection{torch.compile}

Compiling a model with\code{torch.compile}in PyTorch 2.0 results in a 2x speedup from 1.5$\sim$.

\begin{lstlisting}[language=Python]
model = GPT(config)
model = torch.compile(model)  # compilation
\end{lstlisting}

\subsection{profiling}

\begin{lstlisting}[language=Python]
with torch.profiler.profile(activities=[torch.profiler.ProfilerActivity.CPU,
                                         torch.profiler.ProfilerActivity.CUDA]) as prof:
    # learning code
prof.export_chrome_trace("trace.json")
\end{lstlisting}

Open with\code{chrome://tracing}in Chrome browser to analyze bottleneck.
